\section{Sequential decisions: complete proofs}
\label{sec:annex-p17-sequential-decisions}

\begin{annexitem}{annex:p17.sequential.policy.law}{thm:p17-sequential-policy-law}{Constructing the law and its finite Markov projection}
\begin{proof}
For \(N=0\), the initial masses are \(\rho(s_0)\), which
sum to one. Suppose the stated masses through \(X_N\)
have total one. For each such history, adjoining
\(A_N=a\) and \(X_{N+1}=s'\) multiplies its mass by
\(\pi_N(a\mid h_N)P_N(s'\mid s_N,a)\).
Summing first over \(s'\), and then over feasible \(a\),
multiplies that history's mass by one. Thus all the
finite-dimensional masses are normalized and consistent.
Intermediate prefixes ending with an action use the same
formula without the last transition and are also consistent.
For a fixed history as starting prefix, the identical
argument normalizes its continuation masses, whether or
not its original mass was positive.

Encode the finite state and common action alphabets by
distinct real numbers. The preceding masses give consistent
Borel laws for every finite prefix of the alternating
real-coordinate sequence \(X_0,A_0,X_1,A_1,\ldots\).
\Cref{thm:p15-kolmogorov-countable} supplies their unique
countably additive probability law on the coordinate
product sigma algebra. Each coordinate lies in its
specified finite alphabet with probability one.
The event that all coordinates have the prescribed
alphabets and every action is feasible is a countable
intersection of such probability-one cylinder conditions,
so has probability one. Restricting to this event
gives the required path law. Its cylinder values also
prove uniqueness. This is the exact real-coordinate
extension interface; no general-state policy extension
theorem has been assumed.

For the final assertion fix \(\rho,\pi\), and let
\(m_t(s,a)=P_\rho^\pi(X_t=s,A_t=a)\),
\(p_t(s)=\sum_a m_t(s,a)\).
If \(p_t(s)>0\), put
\(\zeta_t(a\mid s)=m_t(s,a)/p_t(s)\).
If \(p_t(s)=0\), choose any probability on \(A_t(s)\).
This is a randomized Markov policy. Its state distribution
at time zero is the same \(\rho\).
If its state distribution at time \(t\) is \(p_t\),
its state/action mass is \(p_t(s)\zeta_t(a\mid s)=m_t(s,a)\),
including a zero-mass state where both sides vanish.
Its next-state mass is therefore
\[
                 \sum_{s,a}m_t(s,a)P_t(s'\mid s,a),
\]
which is also the next-state mass under \(\pi\).
Induction proves equality of all the stated marginals.
The state/action/next-state triples agree as well,
since both use the same controlled transition law.
Thus additive finite-stage reward expectations agree.
The construction uses the marginals under the chosen
\(\rho\); it does not assert one such projection valid
for all initial distributions, or equality of full
history laws.
\end{proof}
\end{annexitem}

\begin{annexitem}{annex:p17.sequential.backward}{thm:p17-sequential-backward}{Backward induction for every history and randomized action}
\begin{proof}
All the state and action sets are finite and nonempty,
and every input reward is finite. Starting with the
finite vector \(g\), the recursion therefore defines
finite vectors and attained maxima at each step.
Choosing the first maximizing action in a fixed order
defines all \(f_t(s)\), including states not reached
from a particular initial distribution.

At time \(N\), every continuation return is \(g(s_N)\).
Suppose the asserted upper bound and attaining policy
have been established at time \(t+1\) for every history.
Under an arbitrary continuation \(\pi\) from \(h_t\),
with last state \(s\), the finite product law gives
\[
 J_{t,N}^\pi(h_t)=
 \sum_a\pi_t(a\mid h_t)\sum_{s'}P_t(s'\mid s,a)
       \left[r_t(s,a,s')
          +\gamma J_{t+1,N}^\pi(h_t,a,s')\right].
\]
The induction hypothesis applies separately to every
extended history \((h_t,a,s')\). Since the probabilities
and \(\gamma\) are nonnegative, replacing its continuation
return by \(v_{t+1}(s')\) gives an upper bound.
The remaining outer sum is a convex combination of
the action values and is at most their maximum \(v_t(s)\).
The deterministic action \(f_t(s)\) attains this maximum,
and the already chosen future maps attain every
continuation value. Thus equality holds for \(f\).
This proves both claims at time \(t\), completing the
backward induction.

The same algebra uses the specified continuation law
at a zero-probability prefix, so the history quantifier
has no hidden positivity restriction.
When \(\gamma=0\), only the immediate reward matters
at a nonterminal time, and the induction remains valid.
When \(N=0\), the initial vector is already \(g\).
Averaging the statewise result against \(\rho\) gives
the stated initial value. Rational additions,
multiplications and comparisons in the finite recursion
produce the exact rational answer; no limiting search
or unspecified optimal continuation remains.
\end{proof}
\end{annexitem}

\begin{annexitem}{annex:p17.sequential.tail}{thm:p17-sequential-tail}{Geometric convergence and the sharp two-policy factor}
\begin{proof}
For \(0\leq\gamma<1\), finite multiplication gives
\[
              (1-\gamma)\sum_{t=0}^{N-1}\gamma^t
                           =1-\gamma^N.
\]
For \(0<\gamma<1\), Bernoulli's inequality, obtained
by multiplying out or by induction, gives
\((1/\gamma)^N\geq1+N(1/\gamma-1)\).
Thus \(\gamma^N\to0\). The case \(\gamma=0\)
is immediate, with the declared zeroth-power convention.
Consequently the geometric series has sum
\(1/(1-\gamma)\), and its tail from \(N\)
has sum \(\gamma^N/(1-\gamma)\).

On every admissible path, \(|r_t|\leq R\).
The absolute reward sum is bounded by this deterministic
geometric series. In particular the partial sums are
uniformly Cauchy functions on path space, and their
limit is a measurable bounded real random variable.
Their expected values converge to its expectation:
the absolute expectation of their difference is at
most their uniform difference.
This proves existence of \(J^\pi\), its norm bound,
and the stated truncation error without interchanging
an uncontrolled infinite sum and an integral.

Two policies that agree through the first \(N\)
actions have identical finite path laws for those
actions and rewards. Their common prefix expectation
cancels when their values are subtracted. Each remaining
tail has absolute expectation at most
\(R\gamma^N/(1-\gamma)\), giving the factor \(2\).
It cannot in general be reduced: take one state
with rewards \(0,R,-R\) available at three actions,
use the zero action for the first \(N\) periods,
and thereafter use respectively \(R\) and \(-R\)
forever. Their difference is exactly
\(2R\gamma^N/(1-\gamma)\).
Finally the added terminal term has absolute expectation
at most \(\gamma^N\|v\|_\infty\).
Adding this bound to the omitted reward tail proves
the last assertion.
\end{proof}
\end{annexitem}

\begin{annexitem}{annex:p17.sequential.contraction}{thm:p17-sequential-contraction}{Completeness, convergence, uniqueness and the monotone criterion}
\begin{proof}
Let \((u_n)\) be Cauchy in supremum norm among bounded
real functions on \(E\).
For each \(x\in E\), the real sequence \(u_n(x)\)
is Cauchy, hence has a real limit \(u(x)\).
Choose \(N\) with \(\|u_n-u_N\|_\infty\leq1\)
for all \(n\geq N\); taking pointwise limits gives
\(|u(x)|\leq\|u_N\|_\infty+1\), so \(u\) is bounded.
For any \(\epsilon>0\), the Cauchy property gives
\(|u_m(x)-u_n(x)|\leq\epsilon\) uniformly in \(x\)
for all sufficiently large \(m,n\).
Letting \(m\) tend to infinity and then taking the
supremum proves \(\|u-u_n\|_\infty\leq\epsilon\).
This proves completeness. A closed subset contains
the limit of every Cauchy sequence of its members.
If each \(u_n\) is measurable for a specified sigma algebra,
its pointwise limit is measurable as well: for every real \(c\),
\[
 \{u>c\}=\bigcup_{m\geq1}\bigcup_{N\geq1}
                    \bigcap_{n\geq N}\{u_n>c+1/m\}.
\]
For the forward inclusion choose \(1/m\) smaller than half
the positive difference \(u(x)-c\); for the reverse inclusion
take the limit of the eventual inequalities.
Thus bounded measurable functions are a closed subspace
of the bounded functions under uniform convergence.

For \(u_{n+1}=Bu_n\), repeated contraction gives
\[
 \|u_{n+1}-u_n\|_\infty
       \leq\gamma^n d,\qquad d=\|Bu_0-u_0\|_\infty .
\]
For \(m>n\), summing these distances gives
\(\|u_m-u_n\|_\infty\leq d\gamma^n/(1-\gamma)\).
The geometric limit just proved makes the sequence
Cauchy; completeness gives a limit \(w\).
The Lipschitz inequality implies continuity of \(B\),
so \(Bw=\lim Bu_n=\lim u_{n+1}=w\).
Taking \(m\to\infty\) in the distance bound gives
the second stated error estimate.
If \(z\) is another fixed point, then
\(\|w-z\|_\infty\leq\gamma\|w-z\|_\infty\);
since \(1-\gamma>0\), the distance is zero.
Finally repeated contraction against this fixed point
gives \(\|u_n-w\|_\infty\leq\gamma^n\|u_0-w\|_\infty\).
If the operator is defined on a nonempty closed
invariant subset, start the iteration there and the
same proof stays in that complete subset.

For the monotone constant-shift criterion, put
\(d=\|u-v\|_\infty\). The pointwise inequalities
\(u\leq v+d\mathbf1\) and \(v\leq u+d\mathbf1\)
give
\[
                   Bu\leq Bv+\gamma d\mathbf1,\qquad
                   Bv\leq Bu+\gamma d\mathbf1.
\]
Together these are exactly the contraction bound.
For \(T_f\), nonnegativity of every transition probability
gives monotonicity and the row sum one gives the
constant-shift identity. A maximum over the finitely
many feasible actions preserves monotonicity, and a
common additive constant commutes with that maximum.
Thus \(T\) has the same identity.
Both operators take finite vectors to finite vectors,
so they are operators on the entire space in question.
\end{proof}
\end{annexitem}

\begin{annexitem}{annex:p17.bellman.bounded.transformations}{thm:p17-bellman-bounded-transformations}{The complete proof of the original bounded equation}
\begin{proof}
Define \((Bu)(p)=\max_k[g_k(p)+h_k(p)u(S_kp)]\)
on the space of all bounded real functions on \(E\).
The condition \(S_kp\in E\) makes every evaluation
well defined. For bounded \(u\),
\(|Bu(p)|\leq C+b\|u\|_\infty\), so \(B\)
maps this space to itself.

For any two finite families of real numbers \(a_k,b_k\),
\[
                       |\max_k a_k-\max_k b_k|
                                      \leq\max_k|a_k-b_k|.
\]
Indeed select an index maximizing the first family;
subtracting the maximum of the second is no larger
than subtracting its entry at that index. Interchange
the two families for the reverse inequality.
Apply this fact with the two expressions defining
\(Bu(p),Bv(p)\). Their difference for each \(k\)
has absolute value at most \(b\|u-v\|_\infty\).
Hence \(B\) is a contraction with modulus \(b\).
The complete fixed-point proof in the preceding
lemma gives existence, uniqueness among all bounded
functions, and convergence from every bounded start.

Starting from the zero function gives nonnegative
iterates because \(g_k,h_k\geq0\); their limit is
nonnegative. At the fixed point,
\(\|F\|_\infty\leq C+b\|F\|_\infty\),
so \(\|F\|_\infty\leq C/(1-b)\).
This proves every assertion, with no compactness,
continuity, finite-state, or measurability assumption
added to the source's bounded-function equation.
The finite deterministic-transition bridge follows
by setting the transition row for action \(k\)
to mass one at \(S_kp\). Its expectation is then
exactly \(u(S_kp)\), as claimed.
\end{proof}
\end{annexitem}

\begin{annexitem}{annex:p17.sequential.stationary.optimum}{thm:p17-sequential-stationary-optimum}{Verification for every randomized history policy}
\begin{proof}
The contraction lemma gives a unique finite vector \(w=Tw\).
For each state, the finite nonempty action set has a
maximizer at \(w\); select one to define \(f^*\).
Then \(T_{f^*}w=w\).
For every feasible action,
\[
               \bar r(s,a)+\gamma\sum_{s'}P(s'\mid s,a)w(s')
                                                        \leq w(s).
\]
At any history under an arbitrary policy \(\pi\),
average this inequality using its action probabilities.
The controlled transition law gives the conditional
inequality for current reward plus discounted \(w(X_{t+1})\).
Equivalently, finite summation over the prescribed
cylinder masses and telescoping yield, from each state \(s\),
\[
 E_s^\pi\left[\sum_{t=0}^{N-1}\gamma^t r_t
                              +\gamma^Nw(X_N)\right]\leq w(s).
\]
This can also be proved directly by induction on \(N\):
the one-step inequality is the base, and the inequality
at the next history cancels its intermediate \(w\) term.
The bound \(\gamma^N\|w\|_\infty\to0\)
and the uniform reward-tail result let \(N\to\infty\),
giving \(J^\pi(s)\leq w(s)\) for every policy.

For \(f^*\), the one-step inequalities are equalities
at every state, so the same finite telescoping identity
and limit give \(J^{f^*}(s)=w(s)\).
This proves \(w=V^*\), attainment, and simultaneous
optimality for all states. Averaging against any
initial distribution preserves the comparison.
It also implies \(\|w\|_\infty\leq R/(1-\gamma)\)
from the policy-value bound.

If \(Tu\leq u\), every action satisfies the same
one-step inequality with \(u\) in place of \(w\).
Finite telescoping and the bounded terminal limit
again give \(J^\pi\leq u\).
For any stationary \(f\), applying the path law
to its first transition and passing to the uniformly
convergent reward series gives \(V^f=T_fV^f\).
If \(TV^f=V^f\), uniqueness of the fixed point of
\(T\) gives \(V^f=w\), so \(f\) is optimal.
If \(f\) is optimal, \(V^f=w\) and its displayed
equation with \(T\) follows.
The same uniqueness argument applies to \(J^\pi\)
for an arbitrary policy when that entire value
vector satisfies \(TJ^\pi=J^\pi\).
Conversely an all-state optimal policy has value
vector \(w\), so satisfies that equation.
\end{proof}
\end{annexitem}

\begin{annexitem}{annex:p17.sequential.value.iteration}{thm:p17-sequential-value-iteration}{Both numerical and policy errors}
\begin{proof}
Contraction against the fixed point gives
\(\|v_k-V^*\|_\infty\leq\gamma^k\|v_0-V^*\|_\infty\).
For \(v_0=0\), the policy-value bound gives the
stated \(R\gamma^k/(1-\gamma)\).
For any \(v\), let \(e=\|v-V^*\|_\infty\).
The triangle inequality and contraction imply
\[
 e\leq\|v-Tv\|_\infty+\|Tv-TV^*\|_\infty
                      \leq\delta(v)+\gamma e.
\]
Rearranging gives the residual estimate.

Put \(L=\|V^*-V^f\|_\infty\).
Because \(T_fv\leq Tv\) by the maximum definition
and \(T_fv\geq Tv-\eta\mathbf1\) by hypothesis,
\(\|Tv-T_fv\|_\infty\leq\eta\).
Use \(TV^*=V^*\), \(T_fV^f=V^f\), and the
triangle inequality to obtain
\[
\begin{split}
 L&\leq\|TV^*-Tv\|_\infty+\|Tv-T_fv\|_\infty
                                    +\|T_fv-T_fV^f\|_\infty\\
  &\leq\gamma e+\eta+\gamma\|v-V^f\|_\infty
   \leq2\gamma e+\eta+\gamma L.
\end{split}
\]
Rearranging and substituting the residual bound
prove the two upper bounds. The statewise
nonnegativity \(V^*-V^f\geq0\) is already proved
by optimality against every policy.
For exact greedy actions, \(T_fv=Tv\), so
\(\eta=0\).

These estimates also supply finite stopping rules.
For value tolerance \(\epsilon>0\), it suffices
that \(\delta(v)\leq(1-\gamma)\epsilon\).
For exact greedy policy loss tolerance \(\epsilon\)
and \(0<\gamma<1\), it suffices that
\(\delta(v)\leq\epsilon(1-\gamma)^2/(2\gamma)\).
Alternatively choose \(k\) by the geometric
prior bound. Such a \(k\) exists because
\(\gamma^k\to0\).
When \(\gamma=0\), \(Tv\) is independent of \(v\)
and an exact greedy action maximizes the sole
relevant reward, proving the separate endpoint.
\end{proof}
\end{annexitem}

\begin{annexitem}{annex:p17.sequential.policy.evaluation}{thm:p17-sequential-policy-evaluation}{The Neumann inverse and a finite exact solve}
\begin{proof}
The identity matrix is row stochastic.
If \(P,Q\) are row stochastic, their product has
nonnegative entries and
\[
                    \sum_j(PQ)_{ij}
                      =\sum_kP_{ik}\sum_jQ_{kj}=1.
\]
Induction proves the assertion for every \(P_f^k\).
For a matrix \(M\), define
\(\|M\|_\infty=\max_i\sum_j|M_{ij}|\).
Finite summation gives
\(\|Mu\|_\infty\leq\|M\|_\infty\|u\|_\infty\)
and
\(\|MN\|_\infty\leq\|M\|_\infty\|N\|_\infty\).
In particular \(\|P_f^k\|_\infty=1\).
The finite geometric bound shows that
\(\sum_{k\geq0}\gamma^kP_f^k\) converges in this norm
to a finite matrix \(B_f\). Its entries are
nonnegative, and summing each row gives
\(1/(1-\gamma)\). Summing a tail gives the
asserted \(\gamma^N/(1-\gamma)\).

For the partial sum \(B_{f,N}=\sum_{k=0}^{N-1}\gamma^kP_f^k\),
finite multiplication telescopes:
\[
       (I-\gamma P_f)B_{f,N}
          =B_{f,N}(I-\gamma P_f)=I-\gamma^NP_f^N .
\]
The right side tends to \(I\) in norm.
Continuity of matrix multiplication, which follows
from the displayed norm inequality, proves
\((I-\gamma P_f)B_f=B_f(I-\gamma P_f)=I\).
This proves the actual two-sided inverse.
The unique solution of \(v=T_fv\) is consequently
\(B_fr_f\). It equals \(J^f\) by the previously
proved first-step equation and uniqueness, or
directly by expanding the finite expected rewards
\(P_f^kr_f\) and passing to the uniform tail limit.

For completeness, an exact finite solve does not
require evaluating an infinite matrix series.
Apply row operations to the augmented matrix
\([\,I-\gamma P_f\mid r_f\,]\).
At column \(j\), earlier pivot columns are unit
columns. Some entry in column \(j\) at or below
row \(j\) must be nonzero. Otherwise assign the
unknown in position \(j\) value one, assign all
later unknowns zero, and solve the earlier pivot
rows for the earlier unknowns. This produces a
nonzero vector in the kernel, contradicting the
proved inverse. Move a nonzero entry to row \(j\),
divide that row by it, and subtract its multiples
from the other rows. These invertible row
operations preserve the solution set and earlier
pivot columns. After \(|S|\) such steps the
coefficient matrix is \(I\), and the last column
is the unique solution. Rational inputs remain
rational throughout. This proves both existence
of every required pivot and finiteness of the
claimed exact algorithm.
\end{proof}
\end{annexitem}

\begin{annexitem}{annex:p17.sequential.improvement}{thm:p17-sequential-improvement}{Howard and Eaton--Zadeh improvement}
\begin{proof}
For any policy \(\pi\), form a new policy which
uses \(g(s)\) initially and, after the successor
state \(s'\) is observed, restarts the original
policy with that state as its initial history.
The specified path kernels and the uniform
discounted tail bound show that its value is
\[
          \sum_{s'}P(s'\mid s,g(s))
                          [r(s,g(s),s')+\gamma J^\pi(s')]
                                                    =T_gJ^\pi(s).
\]
This identifies the source's first-step concatenation.

If \(T_gv\geq v\), monotonicity of \(T_g\) gives
\[
                 v\leq T_gv\leq T_g^2v\leq\cdots .
\]
Contraction makes these iterates converge to
their unique fixed point \(V^g\). Taking limits
proves \(V^g\geq T_gv\geq v\).
For a greedy \(g\) at \(V^f\),
\(T_gV^f=TV^f\geq T_fV^f=V^f\),
so the hypothesis holds.
Subtract \((I-\gamma P_g)V^f\) from the
policy equation for \(V^g\). It gives
\[
                 (I-\gamma P_g)(V^g-V^f)
                              =T_gV^f-V^f=d.
\]
The nonnegative inverse \(B_g\) gives
\(V^g-V^f=B_gd\).
Since \(B_g=I+\sum_{k\geq1}\gamma^kP_g^k\)
and \(d\geq0\), this vector is at least \(d\),
which proves the strictness assertion.

For the second routine put \(v=\max(V^f,V^g)\)
coordinatewise. At a state selecting \(f\),
monotonicity gives
\[
             T_hv(s)=T_fv(s)\geq T_fV^f(s)
                                  =V^f(s)=v(s).
\]
At a state selecting \(g\), the identical
calculation with \(g\) gives \(T_hv(s)\geq v(s)\).
Ties select \(f\) as specified and satisfy the
first calculation. Thus \(T_hv\geq v\).
The monotone iteration argument just proved,
which works for every bounded vector \(v\),
gives \(V^h\geq v=\max(V^f,V^g)\).
This proves both source improvement procedures
with all their comparison and tie conditions.
\end{proof}
\end{annexitem}

\begin{annexitem}{annex:p17.sequential.policy.iteration}{thm:p17-sequential-policy-iteration}{Why the finite algorithm stops}
\begin{proof}
There are exactly \(M=\prod_s|A(s)|\)
deterministic stationary policies.
Evaluation of any one is possible by the
preceding inverse and elimination proof.
For each state, the finite action maximum
exists, so the stipulated improvement is defined.
If it changes no action, each old action is
maximizing at \(V^f\); hence
\(TV^f=T_fV^f=V^f\), which proves optimality.

If it changes an action at state \(s\), the
retain-on-tie rule implies that the old action
was not maximizing there. Thus
\((T_gV^f-V^f)(s)>0\).
The improvement theorem gives \(V^g\geq V^f\)
everywhere and strict inequality at that state.
Along successive nonterminal iterations the
value vectors therefore never decrease at any
coordinate and increase strictly at some
coordinate at each step. A previously visited
policy cannot recur: its identical value vector
would contradict an intervening strict increase
and later coordinatewise nondecrease.
There are only \(M\) policies, so at most
\(M-1\) strict steps occur before stopping.
Rational evaluation and the finite comparisons
are exact and effective for rational data.
The mathematical existence of maxima on a
finite family of arbitrary real numbers
does not, by itself, give a machine a way
to compare unspecified representations of
those real numbers; this explains the
stated algorithmic input boundary.
\end{proof}
\end{annexitem}

\begin{annexitem}{annex:p17.sequential.maintenance}{thm:p17-sequential-maintenance}{Every action value in the worked problem}
\begin{proof}
With \(v_3=(0,0)\), the last-period values are
\[
 v_2(G)=\max(3,1)=3,\qquad
 v_2(B)=\max(0,-2)=0.
\]
At time one, the good-state action values
are
\[
        3+\tfrac12(3)+\tfrac12(0)=\tfrac92
                  \quad\hbox{and}\quad 1+3=4,
\]
so running is optimal. The bad-state values
are \(0+0=0\) and \(-2+3=1\), so repair is
optimal. At time zero, the good-state values
are
\[
          3+\tfrac12(\tfrac92)+\tfrac12(1)=\tfrac{23}4,
                    \qquad 1+\tfrac92=\tfrac{11}2;
\]
the bad-state values are \(1\) and
\(-2+\tfrac92=\tfrac52\).
These comparisons prove the entire table.
They also show exactly why the optimal action
in the bad state changes at the last period:
repair's benefit is a next-state opportunity,
not an additional immediate reward.

For \(\gamma=3/4\), policy
\(f=(\text{run},\text{wait})\) has equations
\[
 V^f(B)=\tfrac34V^f(B),\qquad
 V^f(G)=3+\tfrac38V^f(G)+\tfrac38V^f(B).
\]
They give \(V^f(B)=0,V^f(G)=24/5\).
At this vector, the two good-state action
values are \(24/5\) and
\(1+\tfrac34(24/5)=23/5\);
the bad-state values are \(0\) and
\(-2+\tfrac34(24/5)=8/5\).
Thus one greedy improvement changes only
wait to repair.

For \(g=(\text{run},\text{repair})\), write
its values \(x=V^g(G),y=V^g(B)\). The equations
are
\[
            x=3+\tfrac38x+\tfrac38y,\qquad
            y=-2+\tfrac34x.
\]
Substitution gives \((11/32)x=9/4\), hence
\(x=72/11,y=32/11\).
At this vector, running and repairing
reproduce \(72/11\) and \(32/11\);
service gives \(1+\tfrac34(72/11)=65/11\),
and waiting gives \(\tfrac34(32/11)=24/11\).
The chosen actions are strictly better at
both states. Thus \(TV^g=V^g\), and the
verification theorem gives \(V^g=V^*\).

Starting value iteration at zero gives
\(v_1=(3,0)\).
At \(v_1\), the good-state comparisons are
\(3+\tfrac38\cdot3=33/8\) against
\(1+\tfrac34\cdot3=13/4\).
The bad-state comparisons are zero against
\(-2+\tfrac34\cdot3=1/4\).
This gives \(v_2=(33/8,1/4)\).
The largest absolute reward is \(R=3\).
The general error bound is consequently
\(3(3/4)^k/(1/4)=12(3/4)^k\).
For the greedy policy the multiplier
\(2\gamma/(1-\gamma)\) equals \(6\),
giving \(72(3/4)^k\).
For \(k=23\), the exact integer inequality
\[
          720\cdot3^{23}
              =67783088755440
                <70368744177664=4^{23}
\]
proves that this loss bound is below \(1/10\).
The separate direct action comparisons prove
zero loss for the identified optimal policy,
without interpreting a small generic bound
as an equality certificate.
\end{proof}
\end{annexitem}

\begin{annexitem}{annex:p17.sequential.boundaries}{thm:p17-sequential-boundaries}{Complete counterexamples at the assumptions}
\begin{proof}
For the one-state reward-one model with
\(\gamma=1\), the \(N\)-period return is \(N\),
so the total return is infinite. Its one-step
equation would require \(v=1+v\), impossible
for a finite real value. With reward zero,
the operator is \(Tv=v\), so every scalar is
a fixed point, although the actual reward
sum is zero. For the two-state alternating
example, partial sums from the positive-reward
state alternate between one and zero.
They have no limit. These are precisely
existence, uniqueness and reward-definition
failures once strict discounting is removed.

For the countable-state unbounded-reward
example, starting at zero gives \(X_t=t\).
The discounted reward at every time is
\((1/2)^t2^t=1\), so partial sums diverge
to infinity. Each transition and individual
reward is well defined, but there is no
uniform \(R\). This cannot occur in a
finite homogeneous table of finite real
rewards, where taking their finite maximum
automatically supplies \(R\).

In the countable-action example every
reward is less than one and nonnegative.
Hence all policy values are at most
\(1/(1-\gamma)\). The policy always using
action \(a=N\) has value
\[
                    \frac{1-1/N}{1-\gamma},
\]
which tends to that upper bound.
Under any randomized history policy the
first action \(A_0\) is a positive finite
integer almost surely. Some integer \(j\)
has \(P(A_0=j)>0\), since these probabilities
sum to one. Therefore
\(E(1/A_0)\geq P(A_0=j)/j>0\).
Already the first reward has that strictly
positive deficit from one, and all later
deficits are nonnegative. Thus
\[
                  J^\pi\leq\frac1{1-\gamma}-E(1/A_0)
                                      <\frac1{1-\gamma}.
\]
This proves nonattainment even for arbitrary
randomized history policies and also when
\(\gamma=0\). It is Blackwell's Example~1
at its actual one-state, countably infinite
action domain.

For the history-information example,
take initial observed state \(X_0=H\),
a fair bit, and only one dummy action
at time zero, with zero reward.
Both initial states deterministically
lead to \(X_1=*\). At time one the only
reward is \(\mathbf1_{\{A_1=H\}}\).
A history policy remembers \(X_0\) and
sets \(A_1=H\), receiving one surely.
A Markov policy at \(X_1=*\) has fixed
probabilities \(p,1-p\) for actions zero
and one, independently of the missing
bit under the stipulated information
restriction. Its expected reward is
\(\tfrac12p+\tfrac12(1-p)=1/2\).
The history-dependent reward cannot be
written as the stipulated function of
\((X_1,A_1,X_2)\) with a bit-independent
terminal state \(X_2\). Thus the recorded
state violates an actual model premise.
Using \((H,*)\) as the time-one state
retains the observed bit and restores
the correct Markov reward and optimal
decision in this example.

In the time-varying one-state example,
the policy choosing \(A_t=t\bmod2\)
earns one every period, for value
\(1/(1-1/2)=2\), which is an upper
bound on every policy. A stationary
randomized policy chooses zero with
some fixed probability \(p\in[0,1]\).
Its expected reward is \(p\) at even
times and \(1-p\) at odd times. Its
discounted value is
\[
             \frac{p+(1/2)(1-p)}{1-(1/2)^2}
                    =\frac{2(1+p)}3\leq\frac43.
\]
Here every state, action, reward and
transition is bounded and finite;
time homogeneity is the missing premise.

Finally, with the two identical zero
actions, both stationary value vectors
are zero and both actions maximize
every evaluation. A rule which changes
to the other tied action at each step
never stops if stopping requires an
unchanged policy. Retaining a current
maximizer makes the prescribed algorithm
stop immediately, at an optimal policy.
\end{proof}
\end{annexitem}

\begin{annexitem}{annex:p17.sequential.source.bridge}{thm:p17-sequential-source-bridge}{Source timing, feasible actions and the exact operator equivalence}
\begin{proof}
Blackwell's first decision is indexed by
one and our first decision by zero.
His reward at decision \(n\) is
\(r(s_n,a_n,s_{n+1})\), weighted by
\(\beta^{n-1}\); replacing \(n-1\) by
\(t\) and \(\beta\) by \(\gamma\)
gives exactly our reward sequence.
His stationary operator \(fq(r+\beta u)\)
is, on finite states, the sum
\[
        \sum_{s'}P(s'\mid s,f(s))
                      [r(s,f(s),s')+\gamma u(s')],
\]
which is \(T_fu\).
To identify his generated supremum
operator, enumerate the common finite
actions and form the sequence of
constant state maps selecting them,
repeating the enumeration thereafter.
A generated map may choose any of
these actions at each state; on the
finite state set the corresponding
partition is measurable.
Thus its generated action choices
are exactly all available actions,
and its supremum operator \(U\)
is our finite maximum \(T\).
His nonrandomized Markov plans are
sequences of state maps; our deterministic
Markov witnesses have exactly that form.
The larger randomized history class
in his definition is included in
the policy comparison actually proved.

For state-dependent feasible actions,
let \(A=\bigcup_sA(s)\), select a default
\(a_s^0\in A(s)\), and define
\[
 \kappa_s(a)=
 \begin{cases}a,&a\in A(s),\\a_s^0,&a\notin A(s).\end{cases}
\]
In the common-action model set
\(P'(s'\mid s,a)=P(s'\mid s,\kappa_s(a))\)
and
\(r'(s,a,s')=r(s,\kappa_s(a),s')\).
Each original policy is a policy in this
model using only feasible actions, so
its laws and values are retained.

For the reverse direction, start with
any policy in the common-action model.
Erase an action \(a\) at state \(s\)
to its effective action \(\kappa_s(a)\).
To account for any information stored
in the erased action names, choose an
initial law \(\rho\) positive at every
state, and define a policy on each
positive-probability erased history
by the conditional distribution of the
next effective action given that entire
erased history. These are ratios of
finite sums of cylinder masses. At a
zero-probability erased history choose
any feasible action.
Given an erased history and effective
action, the next-state distribution
is the original \(P\): all unerased
histories and action names being
averaged have that same transition
row by the definition of \(P'\).
Induction on the history length
therefore proves that the new policy
reproduces the erased cylinder law.
Rewards are identical path functions
under the erasure. Because the initial
state is retained in every erased
history and \(\rho(s)>0\), conditioning
this equality on \(X_0=s\) proves it
from every initial state, not merely
after averaging against \(\rho\).
Uniform reward tails extend equality
to infinite discounted values.
Thus the common-action source convention
and the feasible-action model have
the claimed exact policy-law bridge.

At a fixed state, the operator for
action \(a\) is
\[
                 L_a(v)=\bar r(s,a)
                              +\gamma\sum_jP(j\mid s,a)v(j).
\]
Equality \(L_a(v)=L_b(v)\) for all \(v\)
at \(v=0\) forces equality of expected
rewards. For \(\gamma>0\), taking
\(v=\mathbf1_{\{j\}}\) then forces
equality of every transition entry.
These conditions plainly suffice.
For \(\gamma=0\), the continuation
term vanishes and expected reward
equality alone suffices and is necessary.
To demonstrate why equality of
successor-specific reward functions
is stronger, take two successor states
of probability \(1/2\) each. Rewards
\((0,2)\) for one action and \((1,1)\)
for another have the same expected
reward one and the same transition
row, hence the same operator for
every continuation vector, although
the realized reward functions differ.
This supplies the exact operator
equivalence used in the finite
construction.
\end{proof}
\end{annexitem}
